% Problem record op_a11fdfadc84e4379. This TeX file is derived from
% database/problems_json/op_a11fdfadc84e4379.json by scripts/sync-tex.mjs; edit the JSON record
% and rerun the script rather than editing this file.
\section{Aaronson–Ambainis conjecture}
\label{sec:op_a11fdfadc84e4379}

\paragraph{Problem.}
Must every bounded low-degree real polynomial on the Boolean cube with
non-negligible variance have a variable of inverse-polynomial influence?

Let $p:\mathbb R^N\to\mathbb R$ be a real multilinear polynomial of degree
at most $d$ satisfying
\begin{equation}
  0\leq p(X)\leq 1
  \qquad
  \text{for every }X\in\{0,1\}^N .
  \label{eq:a11f-bounded}
\end{equation}
For $i\in[N]$, let $X^{(i)}$ denote $X$ with its $i$-th bit flipped, and
define the influence of the $i$-th variable by
\begin{equation}
  \operatorname{Inf}_i[p]
  =
  \mathbb E_{X\in\{0,1\}^N}
  \left[
    \left(
      p(X)-p(X^{(i)})
    \right)^2
  \right].
  \label{eq:a11f-influence}
\end{equation}
Define the variance of $p$ on the uniform Boolean cube by
\begin{equation}
  \operatorname{Var}[p]
  =
  \mathbb E_{X\in\{0,1\}^N}
  \left[
    \left(
      p(X)-\mathbb E[p]
    \right)^2
  \right].
  \label{eq:a11f-variance}
\end{equation}
Suppose that for some $\varepsilon>0$,
\begin{equation}
  \operatorname{Var}[p]\geq\varepsilon .
  \label{eq:a11f-nontrivial}
\end{equation}
The Aaronson--Ambainis conjecture asks whether there is a universal constant
$C>0$ such that every polynomial satisfying \eqref{eq:a11f-bounded} and
\eqref{eq:a11f-nontrivial} has some coordinate $i\in[N]$ whose influence,
as defined in \eqref{eq:a11f-influence}, satisfies
\begin{equation}
  \operatorname{Inf}_i[p]
  \geq
  \left(\frac{\varepsilon}{d}\right)^C .
  \label{eq:a11f-conjecture}
\end{equation}
Equivalently, the conjecture asks whether the maximum influence can always be
bounded below by a fixed polynomial in $1/d$ and the variance
\eqref{eq:a11f-variance}, independently of the ambient dimension $N$.

\subsection*{Status}

Unsolved

\subsection*{Source}

Explicitly proposed by Scott Aaronson and Andris Ambainis in Aaronson and
Ambainis (2014), Conjecture~1.7, under the name ``Bounded Polynomials Have
Influential Variables''
\sourcecite{ref:a11f-aa}{AA14}.  The related statement asserting
polynomial-query classical simulation of a quantum query algorithm on most
inputs appears separately as Conjecture~1.5 in the same paper and is
explicitly described there as folklore; it should therefore not be
attributed as an original Aaronson--Ambainis conjecture.  The formulation in
the statement above is a self-contained restatement of Conjecture~1.7 rather
than a verbatim quotation.

\subsection*{Progress}

\begin{itemize}
  \item Beals, Buhrman, Cleve, Mosca, and de Wolf showed that the acceptance
probability of a quantum algorithm making $T$ oracle queries is represented
by a real multilinear polynomial of degree at most $2T$; bounded low-degree
polynomials such as those in the statement therefore capture a fundamental
aspect of quantum query algorithms
\sourcecite{ref:a11f-bbc}{BBCMW01}.

  \item Aaronson and Ambainis proposed the influence conjecture under the name
``Bounded Polynomials Have Influential Variables'' and showed it would imply
a strong classical-simulation principle: for every $T$-query quantum
algorithm and every positive approximation and failure parameters, a
deterministic classical algorithm using a polynomial number of queries can
approximate the quantum algorithm's acceptance probability on all but a
small fraction of Boolean inputs.  The conjecture would thereby formalize
that superpolynomial quantum query speedups require substantial promise
structure rather than occurring generically on the Boolean cube
\sourcecite{ref:a11f-aa}{AA14}.

  \item The influential-variable conjecture should be distinguished from the
simulation statement: in the same paper the almost-everywhere
classical-simulation statement appears separately as a folklore conjecture,
while the influence conjecture is the route Aaronson and Ambainis introduced
toward it
\sourcecite{ref:a11f-aa}{AA14}.

  \item Bansal, Sinha, and de Wolf proved the desired inverse-polynomial influence
bound for completely bounded degree-$d$ block-multilinear forms with
constant variance, giving almost-everywhere classical simulation for a
corresponding class of quantum query algorithms including iterated
Forrelation; this does not establish the conjecture for arbitrary bounded
low-degree polynomials
\sourcecite{ref:a11f-bsw}{BSW22}.

  \item Bhattacharya proved a partial result via random restrictions: a bounded
degree-$d$ polynomial becomes well approximated by a polynomial-size junta
after an appropriate random restriction with high probability, and the
Aaronson--Ambainis conclusion holds for a non-negligible fraction of such
restrictions under an appropriate variance condition.  This establishes
additional structure in bounded low-degree polynomials but does not prove
the conjecture without restriction
\sourcecite{ref:a11f-bhatt}{Bhat25}.

  \item Blanc, Docter, Strassle, and Tan proved the associated almost-everywhere
classical-simulation conjecture for bounded-round quantum query algorithms:
a $t$-query, $d$-round quantum algorithm can be simulated on most inputs
using $t^{O(d^2)}$ classical queries, giving polynomial overhead for
constant $d$.  Their result does not prove the influence conjecture for
arbitrary bounded polynomials and does not settle the unrestricted adaptive
quantum-query simulation problem
\sourcecite{ref:a11f-bdst}{BDST26}.
\end{itemize}

\subsection*{References}

\begin{enumerate}[label={},leftmargin=4.8em,labelsep=0.5em]
  \footnotesize
  \item[\textup{[AA14]}]\label{ref:a11f-aa}
  Scott Aaronson and Andris Ambainis, ``The Need for Structure in Quantum
Speedups,'' \emph{Theory of Computing} \textbf{10}, 133--166 (2014).
\href{https://doi.org/10.4086/toc.2014.v010a006}{doi:10.4086/toc.2014.v010a006};
\href{https://arxiv.org/abs/0911.0996}{arXiv:0911.0996}.

  \item[\textup{[BBCMW01]}]\label{ref:a11f-bbc}
  Robert Beals, Harry Buhrman, Richard Cleve, Michele Mosca, and Ronald de
Wolf, ``Quantum Lower Bounds by Polynomials,'' \emph{Journal of the ACM}
\textbf{48}, 778--797 (2001).
\href{https://doi.org/10.1145/502090.502097}{doi:10.1145/502090.502097};
\href{https://arxiv.org/abs/quant-ph/9802049}{arXiv:quant-ph/9802049}.

  \item[\textup{[BSW22]}]\label{ref:a11f-bsw}
  Nikhil Bansal, Makrand Sinha, and Ronald de Wolf, ``Influence in
Completely Bounded Block-Multilinear Forms and Classical Simulation of
Quantum Algorithms,'' in \emph{37th Computational Complexity Conference
(CCC 2022)}, LIPIcs \textbf{234}, 28:1--28:21 (2022).
\href{https://doi.org/10.4230/LIPIcs.CCC.2022.28}{doi:10.4230/LIPIcs.CCC.2022.28};
\href{https://arxiv.org/abs/2203.00212}{arXiv:2203.00212}.

  \item[\textup{[Bhat25]}]\label{ref:a11f-bhatt}
  Sreejata Kishor Bhattacharya, ``Random Restrictions of Bounded Low Degree
Polynomials Are Juntas,'' in \emph{16th Innovations in Theoretical Computer
Science Conference (ITCS 2025)}, LIPIcs \textbf{325}, 17:1--17:21 (2025).
\href{https://doi.org/10.4230/LIPIcs.ITCS.2025.17}{doi:10.4230/LIPIcs.ITCS.2025.17};
\href{https://arxiv.org/abs/2402.13952}{arXiv:2402.13952}.

  \item[\textup{[BDST26]}]\label{ref:a11f-bdst}
  Guy Blanc, Jordan Docter, Carmen Strassle, and Li-Yang Tan, ``Quantum
Speedups Require Structure or Depth,'' to appear in \emph{67th IEEE Symposium
on Foundations of Computer Science (FOCS 2026)} (2026).
\href{https://arxiv.org/abs/2608.19158}{arXiv:2608.19158}.
\end{enumerate}

\subsection*{Comment}

The central unresolved question is whether the dimension-independent
influence lower bound in the statement holds for every bounded low-degree
polynomial, without additional multilinearity structure, random restriction,
or bounded-adaptivity assumptions on the associated quantum query
algorithms.

\subsection*{Field}

Quantum algorithm

\subsection*{Topic}

Computational complexity and computability

\subsection*{ID}

\texttt{op\_a11fdfadc84e4379}
